\documentclass{article}
%\usepackage{tgpagella}
%\usepackage{tgtermes}
\usepackage[T1]{fontenc}
\usepackage{tgtermes} 
\usepackage{siunitx}
\usepackage{newtxmath} 
\usepackage{chemfig}
\usepackage{chemformula}

\begin{document}
Pause Hier ist ein blöder Text $\SI{1}{\joule}$ und Oskar stinkt!"\\
Das ist ein zufälliger deutscher Text zur Demonstration. Wörter wie Haus, Baum, Auto, Sonne, Mond und Sterne  $x=m^2$
erscheinen hier in unerwarteten Kombinationen. Die 
Sätze sind kurz und bündig, aber nicht immer sinnvoll. 
Manchmal fehlt sogar die Verbindung zwischen den Gedanken.
$$m=c^2\frac{dIl}{mnop}$$

%{tgtermes}
klassisch $$\ch{T2 + 100 AAA2O5 -> 2 + M^3+ + e- + ext + C}$$

\end{document}